\documentclass[11pt,a4paper]{article}

% --- Packages ---
\usepackage[utf8]{inputenc}
\usepackage[T1]{fontenc}
\usepackage[french]{babel}
\usepackage[margin=0.8in]{geometry}
\usepackage{xcolor}
\usepackage{hyperref}
\usepackage{enumitem}

% --- Custom Colors ---
\definecolor{primary}{RGB}{0, 51, 102} % Dark Blue

% --- Formatting ---
\pagestyle{empty}
\setlength{\parindent}{0pt}
\setlength{\parskip}{0.6em}

\begin{document}

% ============================================================
% SENDER INFO
% ============================================================
{\bfseries Omar Bouhlal} \\
Meknès, Maroc \\
+212 766 678 608 \\
omarbouhlal05@gmail.com \\
\href{https://linkedin.com/in/omar-bouhlal-101ba6353}{linkedin.com/in/omar-bouhlal-101ba6353} \\
\href{https://github.com/OmarBouhlal}{github.com/OmarBouhlal}

\vspace{1.5em}

% ============================================================
% RECIPIENT INFO (GENERIC)
% ============================================================
\hfill {\bfseries À l'attention de l'Établissement d'Accueil} \\
\hfill Candidature à la Mobilité Académique Internationale

\vspace{1em}

% ============================================================
% DATE & SUBJECT
% ============================================================
Fait à Meknès, le \today

\vspace{1em}

{\bfseries Objet : Candidature pour un programme de mobilité internationale / Double Diplôme}

\vspace{1.5em}

% ============================================================
% LETTER BODY
% ============================================================
Madame, Monsieur,

Actuellement étudiant en 4ème année (2ème année du cycle ingénieur) à l'École Nationale Supérieure d'Arts et Métiers (ENSAM) de Meknès, je vous adresse ma candidature pour intégrer un programme de mobilité internationale dans le cadre de mon cursus académique pour l'année [ANNÉE].

Mon parcours à l'ENSAM Meknès, et plus particulièrement ma formation en Génie Logiciel et Systèmes Intelligents, m'a permis de développer une solide rigueur scientifique. Cette discipline m'a confronté à la conception de systèmes complexes, exigeant une structuration de la pensée et un entraînement constant de la logique. Au-delà de l'aspect technique, cette formation a forgé ma capacité à résoudre des problèmes multidisciplinaires avec méthode et précision, des compétences que je considère comme fondamentales pour tout ingénieur.

L'opportunité de poursuivre une partie de mes études à l'étranger représente une étape cruciale dans mon projet professionnel. Je suis particulièrement motivé par la perspective d'évoluer dans un environnement académique de haut niveau, tout en découvrant de nouvelles méthodes de travail et d'apprentissage. Mon objectif est d'élargir mes horizons et de bénéficier d'une ouverture internationale indispensable pour ma future carrière, quel que soit le domaine d'application final.

Au-delà de l'aspect académique, je suis convaincu qu'une expérience internationale au sein d'un environnement multiculturel me permettra de développer une plus grande adaptabilité et d'enrichir ma perspective sur les enjeux technologiques mondiaux. Ma curiosité intellectuelle et ma capacité à travailler en équipe sont des atouts que je souhaite mettre au service de l'établissement d'accueil.

Je reste à votre entière disposition pour tout entretien ou information complémentaire. Je vous joins mon curriculum vitae pour plus de détails sur mon parcours.

Je vous prie d'agréer, Madame, Monsieur, l'expression de mes salutations distinguées.

\vspace{2.5em}

% ============================================================
% SIGNATURE
% ============================================================
Omar Bouhlal

\end{document}
